\section{讨论}

本文首先提示了一个值得进一步展开的 scaling law 方向。在当前结果中，随着验证数据规模从 $10^3$ 量级扩展到 $2\times10^4$ 量级，并配合更强的 surrogate 与生成模型，整体设计质量呈现持续改善趋势。从这个意义上说，本文其实已经给出了一个初步的、属于本问题本身的 scaling-law 信号。它与大语言模型中经验 scaling law 的思路有相通之处，但并不完全等同\cite{kaplan2020scaling}：当目标是结构化的 wavelength--theta--polarization 响应场时，模型能力不仅取决于生成器本身，也取决于高价值响应模式的数据覆盖，以及前向代理对这些模式的分辨能力。当前的局限在于，现有 scaling 主要还发生在数据量和模型能力上，而材料种类、周期大小以及相位目标等物理自由度尚未被共同纳入统一框架。未来更有价值的方向，是把这些物理设计轴也一起扩展，并配合更密的电磁验证，从而把这套方法逐步推向更通用的超表面逆向设计。

其次，本文的方法并不局限于文中展示的几类任务。由于目标被写成统一的响应张量，这套框架可以自然扩展到偏振--频率--角度三重联合设计，让多个物理维度同时参与调控。这样的设计语言与紧凑型折叠光谱仪和角分辨光谱仪\cite{faraji2018folded,cai2024angleresolved}、多维传感与计算型探测\cite{jiang2024multidimensional}、以及偏振编码信息与光学加密器件\cite{guo2022stokes}等方向都高度相关。从这个意义上说，本文更重要的价值不只是给出几个具体器件，而是提供了一种可以继续外推的响应级设计框架。
